% =====================================================================
%  REGISTRO DELLE MODIFICHE di analisi-dei-capitolati.tex
%
%  Le righe le aggiunge il bot dopo il merge di una pull request con
%  un'etichetta di versione (correzione, contenuto, approvazione):
%  GUIDA.md, sezione 5. La riga più recente è in cima.
%  Si può correggere a mano (per esempio un refuso in una descrizione)
%  con una pull request SENZA etichetta. La classe legge questo file da
%  sola: nel .tex non serve scrivere nulla.
% =====================================================================
\modifica{1.0.0}{2026-10-28}{spagnoli}{-}{Approvazione del documento}
\modifica{0.10.0}{2026-10-28}{guglielmo}{corro}{Termini marcati allineati al Glossario}
\modifica{0.9.0}{2026-10-27}{guglielmo}{cescobolla}{Revisione complessiva: sintesi dei contenuti e approfondimento dei capitolati C3--C9}
\modifica{0.8.0}{2026-10-27}{biasini}{corro}{Integrazione degli esiti dell'incontro con Eggon del 26 ottobre}
\modifica{0.7.0}{2026-10-26}{guglielmo}{cescobolla}{Valutazione finale di C2, tabelle di confronto e conclusioni}
\modifica{0.6.0}{2026-10-25}{biasini}{corro}{Stesura della sintesi dei capitolati C3--C9}
\modifica{0.5.0}{2026-10-24}{guglielmo}{cescobolla}{Stesura dell'analisi del capitolato scelto, C2 -- NEXUM Concierge}
\modifica{0.4.0}{2026-10-23}{biasini}{corro}{Stesura dell'analisi di C1 -- AURA con gli esiti dell'incontro del 22 ottobre}
\modifica{0.3.0}{2026-10-22}{guglielmo}{cescobolla}{Stesura dell'analisi di C0 -- AIRA con gli esiti dell'incontro del 21 ottobre}
\modifica{0.2.0}{2026-10-21}{biasini}{corro}{Stesura della metodologia e dei criteri di valutazione}
\modifica{0.1.0}{2026-10-20}{guglielmo}{cescobolla}{Prima stesura: struttura del documento e introduzione}
